\subsubsection{গন্তব্য বোঝা যায় এমন লিংক: ছোট অনুশীলন}
শিক্ষার্থী পরীক্ষার রুটিন খুঁজছে। ``এখানে ক্লিক'' দেখে গন্তব্য বোঝা কঠিন; ``পরীক্ষার রুটিন'' নামে লিংক দিলে সিদ্ধান্ত সহজ হয়। নিচের কোডটি পেজের \texttt{body}-তে লেখা হবে।
\par\noindent\begin{minipage}[t]{.50\linewidth}\vspace{0pt}\small
\textbf{\color{chFourTeal}কোড}\par
\begin{verbatim}
<a href="exam-routine.html">
  Examination routine
</a>
\end{verbatim}
\end{minipage}\hfill
\begin{minipage}[t]{.46\linewidth}\vspace{0pt}\small
\textbf{\color{chFourTeal}আগে ভাবো, পরে পরীক্ষা করো}\par
দৃশ্যমান লেখা Admission notice করলেও গন্তব্য বদলায় না; গন্তব্য থাকে \texttt{href}-এ। লেখাটি গন্তব্যের সঙ্গে না মিললে ব্যবহারকারীর ভুল প্রত্যাশা তৈরি হয়। লিংকের ঠিকানা ও দৃশ্যমান পরিচয় দুটোই মিলতে হয়।
\end{minipage}\par
